\section{OpenAI、GPT-6とIntelligent UIをChatGPTへ展開}
OpenAIはGPT-6とIntelligent UIのChatGPT向けグローバル展開を公式に発表した。投稿ではPlus / Pro / Business / Enterpriseを当日、Free / Goを翌日から順次対象とし、有料側はGPT-6 Sol、Free / Go側はGPT-6 Lunaを利用すると説明する。WorkとCodex向けモデルは今回のリリースでは変更しないとしている。公式ブログURLはGrok収集では未取得である。
\dailyxsource{OpenAI (@OpenAI)}{https://x.com/OpenAI/status/2107894997538525580}

\section{Anthropic、Claude Haiku 5.5を発表}
AnthropicはClaude Haiku 5.5を発表し、Haiku 4.5比で平均約75\%低コストと主張した。公式ページでは入力・出力価格、長文時の別価格、Sonnet 5.5のcache reads半額化、Python / TypeScript SDKでのcomputer use・browser useベータも案内している。ベンチマーク値はAnthropic自身の評価で、独立再測定はGrok収集では行われていない。
\dailyxsource{Claude (@claudeai)}{https://x.com/claudeai/status/2107894039626277339}

\section{OpenAI、ChatGPT for TeensとCollege Plannerを案内}
OpenAIは18歳未満のアカウント向けChatGPT for Teensの進捗を共有し、保護機能をデフォルトで有効にすると投稿した。あわせて、米国の四年制大学進学を計画する高校生向けCollege Plannerをcoming soonとしてプレビューし、要件、締切、タスク、学資援助手順をまとめる構想を示した。詳細な公式ポリシー文書は今回未確認である。
\dailyxsource{OpenAI (@OpenAI)}{https://x.com/OpenAI/status/2107909792945832070}

\section{Google、SynthIDの公開検証ポータルを公開}
Google AIはSynthIDウォーターマークを一般利用者が検証できるsynthid.comを案内した。投稿ではGoogleに加えOpenAI、NVIDIA、Kakaoなどのパートナー生成物も対象とし、画像・動画1800億件、音声24万年以上分へのウォーターマーク付与、日次約100万件の検証リクエストという規模を示している。これらの数値はGoogle側の主張である。
\dailyxsource{Google AI (@GoogleAI)}{https://x.com/GoogleAI/status/2107836750257107133}

\section{Isomorphic Labs、Virtual Biology Initiativeへ参加}
Isomorphic LabsはVirtual Biology Initiativeの創設メンバーとして参加すると発表した。公式投稿は、疾患理解と治療開発を加速するためのオープンリソース構築への貢献を掲げ、Demis HassabisもCZIとの連携と医療・健康へのAI適用として言及した。投稿内の記事URLは確認されているが、Grok収集では記事本文の詳細スコープまでは確認していない。
\dailyxsource{Isomorphic Labs (@IsomorphicLabs)}{https://x.com/IsomorphicLabs/status/2107820493940122052}

\section{Grok Bot、Xの検索・閲覧・監視機能を告知}
Grok Bot公式アカウントは、X上の情報をsearch、read、monitorできるようになったと投稿した。Elon Muskは多くのBotリクエストを将来の高速Grok 4.8で処理する計画にも言及したが、4.8は観測時点で未リリースの将来事項である。第三者による「1プロンプトで100件超の投稿分析」などの主張は公式本文には含まれない。
\dailyxsource{Grok Bot (@bot)}{https://x.com/bot/status/2107949161878606089}
